BACKGROUND. Epithelial-mesenchymal plasticity (EMP) drives metastasis and therapy resistance, and hybrid E/M cells carry higher metastatic potential than cells at either extreme. The field scores plasticity as a single quantity, and every established metric, from EMT scores to heterogeneity indices, measures a distribution of cells at one moment. Such a measurement cannot distinguish a population of stable cells with varied specializations from a population whose cells convert between states, yet these should not seed metastases equally, since metastasis imposes conflicting demands in sequence and clusters of circulating tumor cells are the efficient seeding unit.
OBJECTIVE AND HYPOTHESIS. We hypothesize that EMP comprises distinguishable dynamical regimes differing in two measurable quantities, reach, the states a cell can access after perturbation, and retention, the transitions it makes within the hybrid range before leaving it; and that the regimes a cell line can access predict its metastatic competence better than its hybrid fraction does.
SPECIFIC AIMS. Aim 1 defines the regimes on multilevel EMT network models and fixes, in advance of any experiment, the signatures each regime predicts, the observation window an assignment requires, and the markers it is read on; a selection layer then applies the stresses of transit and colonization to simulated trajectories and yields a predicted ordering of regimes by survival. Aim 2 tests the regimes in roughly 30 mouse and human breast cancer lines and against metastatic outcome. Aim 3 asks which molecular mechanisms set the graded regulation that produces these regimes. Aims 1 and 2 together test the hypothesis; Aim 3 builds on them.
STUDY DESIGN. Reach and retention are computed on state-transition graphs across a noise sweep at matched state-space size, so that each network is described by the regimes it can access rather than by behavior at one noise level. Predictions are registered before experimental results are examined. In cells, regime assignment uses existing mass cytometry and flow data from the collaborating laboratory, in which functional markers of stemness, apoptotic state, hypoxia and cell cycle are measured alongside epithelial and mesenchymal markers in single cells; sorted populations followed to re-equilibration separate a mixture of committed cells from a converting population, and a TGF-beta dose series measures response steepness per line. Fourteen of these lines have been through orthotopic and tail vein metastasis experiments, so regime assignment is tested against those outcomes with hybrid fraction entered as the competing predictor; comparing the two entry points separates effects in transit from effects at escape. Mechanisms are tested by scoring read-level methylation disorder against regime assignment, and by predicting the behavior of existing inducible miR-200 and ZEB1 lines. Each aim states the result that would refute it.